\section{模型假设与歧义处理}
缓存大小沿用题面的抽象容量单位，不将其转换成字节；执行时间统一使用时钟周期。缓存申请与释放耗时为零，操作耗时固定为节点的周期字段。原始依赖边均不可删除，也不依据真实芯片的额外机制添加题面没有给出的吞吐假设。

问题一的驻留量按所有缓存类型合计计算，而非各类型峰值的最大值。对同一类型同时驻留的缓冲区要求地址不重叠，不同类型拥有独立地址空间。

附录 C 通过前一缓冲区的释放节点约束后一缓冲区的申请节点，直接描述了未发生换出的地址复用。发生换出时，物理地址在换出结束后即可复用，缓冲区的逻辑生命则持续到原始释放节点。本文将两者区分为逻辑生命周期与物理驻留段：申请或换入开启一段驻留，释放或换出结束该段。新驻留段必须等待所复用旧段的结束。若仍以最终释放节点约束换出期间的所有复用，可能人为形成依赖环，无法表达题面允许的换出机制。因此驻留段是本文对该情形的明确补充解释；与未提供的官方评测程序的一致性尚不能确证。

换出目标曾被任一输入复制节点引用时，按题面规定只计换入搬运量，换出时间为零；不自行判断数据是否被后续修改并改变计费方式。所有缓存类型均依题面统一换出规则处理。问题三要求额外搬运量不高于问题二选定方案，从而避免人为设定宽松的容忍阈值。
